\documentclass[11pt]{article}

\usepackage[margin=1in]{geometry}
\usepackage{amsmath, amssymb}
\usepackage{booktabs}
\usepackage{enumitem}
\usepackage{xcolor}
\usepackage{hyperref}
\usepackage{listings}
\usepackage{microtype}

\hypersetup{colorlinks=true, linkcolor=blue, urlcolor=blue}

\lstset{
  basicstyle=\ttfamily\small,
  columns=fullflexible,
  frame=single,
  framesep=4pt,
  rulecolor=\color{gray!40},
  showstringspaces=false,
  breaklines=true
}

\newcommand{\code}[1]{\texttt{#1}}
\newcommand{\Up}{\mathbf{U}_p}
\newcommand{\Lp}{\boldsymbol{\Lambda}_p}
\newcommand{\bmu}{\boldsymbol{\mu}}
\newcommand{\R}{\mathbb{R}}

\title{Feature-Level Anomaly Generation Experiments\\
       \large Experiment Plan, Hypotheses, and Interpretation}
\author{SimpleNet feature-anomaly research project}
\date{}

\begin{document}
\maketitle

\begin{abstract}
This document explains the experiment suite implemented for the SimpleNet
feature-anomaly project. The purpose is to test whether synthetic anomalies at
the feature level can be made more useful than SimpleNet's original random
Gaussian noise. Each experiment changes one controlled aspect of the anomaly
generator: anchor, radius calibration, spatial sparsity, or discriminator-guided
refinement. The intended output is not only an AUROC table, but a causal story
about which geometric property of feature-space anomalies matters for image
detection and pixel localization.
\end{abstract}

\tableofcontents

%======================================================================
\section{Research Question}
%======================================================================

SimpleNet trains a discriminator using normal patch features as positives and
synthetic anomalous features as negatives. In the original setup, fake features
are produced by adding isotropic Gaussian noise:
\begin{equation}
  \mathbf{x}_{\mathrm{fake}}
    = \mathbf{x}_{\mathrm{real}} + \boldsymbol{\epsilon},
  \qquad
  \boldsymbol{\epsilon} \sim \mathcal{N}(0, \sigma^2 I).
\end{equation}

\noindent
This is simple and fast, but it ignores the geometry of the normal feature
distribution. The central hypothesis of this project is:

\begin{quote}
Feature-level anomalies are more useful when they are anchored on real normal
features, move in meaningful local feature directions, have calibrated
magnitude, and provide spatially localized supervision.
\end{quote}

\noindent
The experiments therefore compare random noise against a geometry-aware family
of anomalies derived from the fitted per-patch model
\begin{equation}
  \Sigma_p = \Up \Lp \Up^\top + \varepsilon_p I,
\end{equation}
where $p$ indexes spatial patch location. In the most important setting, an
anomaly is generated as
\begin{equation}
  \mathbf{x}_{\mathrm{fake},p}
    =
  \mathbf{x}_{\mathrm{real},p}
    +
  r_p\, \Up \sqrt{\Lp}\,\mathbf{v},
  \qquad
  \mathbf{v} \sim \mathrm{Uniform}(\mathbb{S}^{k-1}).
\end{equation}

%======================================================================
\section{Provenance of the Ideas}
%======================================================================

This section records where each idea in the experiment suite comes from. The
goal is not to claim that every component is novel in isolation. The goal is to
make clear which prior ideas are being combined and what new question this
project asks.

\subsection{SimpleNet: the Baseline to Improve}

The starting point is SimpleNet~\cite{simplenet}. SimpleNet introduced the
specific training pattern used in this project: extract local features from a
pretrained backbone, optionally adapt them, synthesize anomalous features by
adding Gaussian noise to normal features, and train a binary discriminator. This
project keeps the SimpleNet discriminator framing but questions whether the
feature-level anomaly generator should be isotropic random noise.

\paragraph{Taken idea.}
Feature-space synthetic negatives can train a discriminator for anomaly
detection.

\paragraph{Project modification.}
Replace directionless Gaussian noise with geometry-aware feature perturbations.

\subsection{PaDiM: Patch-Wise Gaussian Feature Geometry}

The use of a per-patch Gaussian model is inspired by PaDiM~\cite{padim}. PaDiM
models normal patch embeddings with multivariate Gaussian distributions and uses
Mahalanobis distance for anomaly scoring. Our project borrows the patch-wise
normal-distribution view, but uses it to synthesize training negatives rather
than only to score test samples.

\paragraph{Taken idea.}
Normal patch features have spatially specific distributions that can be modeled
with Gaussian statistics.

\paragraph{Project modification.}
Use the fitted covariance directions $\Up$ and scales $\Lp$ to generate fake
features for discriminator training.

\subsection{PatchCore: Respecting the Normal Feature Manifold}

PatchCore~\cite{patchcore} motivates the broader view that nominal patch
features themselves are the important object. It uses a representative memory
bank of normal patch features and nearest-neighbor distance at inference. This
project does not use a memory bank for scoring, but the anchored generator is
conceptually aligned with PatchCore's emphasis on real normal patch features:
fake samples should start from real normal anchors rather than from an abstract
mean alone.

\paragraph{Taken idea.}
Real normal patch features provide a strong empirical representation of the
normal manifold.

\paragraph{Project modification.}
Anchor synthetic negatives on real training features:
\[
  \mathbf{x}_{\mathrm{fake},p}
  =
  \mathbf{x}_{\mathrm{real},p} + \Delta_p.
\]

\subsection{CutPaste, NSA, DRAEM, and DeSTSeg: Synthetic Local Supervision}

Image-level anomaly synthesis methods motivate the sparse-mask experiments.
CutPaste~\cite{cutpaste} uses a simple cut-and-paste augmentation as a
self-supervised anomaly signal. NSA~\cite{nsa} creates more natural synthetic
anomalies by blending patches. DRAEM~\cite{draem} and DeSTSeg~\cite{destseg}
show that synthetic corruptions and their masks can provide direct localization
supervision.

\paragraph{Taken idea.}
Synthetic anomalies are more useful for localization when they are spatially
local and come with an implicit or explicit mask.

\paragraph{Project modification.}
Apply synthetic perturbations only to selected feature patches and compute fake
loss only on those patches:
\[
  \mathcal{L}_{\mathrm{fake}}
  =
  \frac{1}{|\Omega|}
  \sum_{p \in \Omega}
  \max(0, D(\mathbf{x}_{\mathrm{fake},p}) + m).
\]

\subsection{GLASS: Gradient-Guided Hard Feature Negatives}

GLASS~\cite{glass} is the closest recent inspiration for the gradient-refinement
experiments. It argues that anomaly synthesis should be controllable and should
cover weak, near-in-distribution defects. Its feature-level synthesis uses
gradient ascent and projection constraints to create useful hard negatives.

\paragraph{Taken idea.}
Feature-level synthetic anomalies should not be purely random; they can be
guided by the current model to produce harder, more informative negatives.

\paragraph{Project modification.}
Start from anchored PCA anomalies, then take discriminator-guided refinement
steps while projecting updates back into the fitted patch PCA subspace and
clamping their Mahalanobis radius.

\subsection{CRAS: Distance-Guided Magnitude Calibration}

CRAS~\cite{cras} motivates the radius-calibration experiments. Its anomaly
synthesis strategy adaptively adjusts noise based on normal-data distribution to
reduce normal/anomaly overlap. This supports our suspicion that a single global
noise scale, or a raw Mahalanobis threshold radius, may be poorly calibrated for
training.

\paragraph{Taken idea.}
Synthetic anomaly magnitude should depend on the position of the normal sample
relative to the normal distribution.

\paragraph{Project modification.}
Introduce \code{patch} and \code{anchor} radius modes:
\[
  r_p = \rho \sqrt{T_p/C}U(0,1),
  \qquad
  r_{i,p}
  =
  \rho
  \max\left(
    \sqrt{T_p} - \sqrt{s(\mathbf{x}_{i,p})},
    0.05\sqrt{T_p}
  \right)U(0,1).
\]

\subsection{Summary of Influence}

\begin{center}
\begin{tabular}{p{0.24\linewidth}p{0.27\linewidth}p{0.37\linewidth}}
\toprule
Project component & Main source idea & What we changed \\
\midrule
Vanilla noise baseline & SimpleNet & Re-run in current controlled path \\
Patch Gaussian/PCA model & PaDiM & Use for synthesis, not only scoring \\
Real-feature anchor & PatchCore & Perturb anchors rather than only score them \\
Sparse patch masks & Image-level anomaly synthesis & Move masks to feature space \\
Gradient refinement & GLASS & Constrain updates to per-patch PCA geometry \\
Radius calibration & CRAS & Use patch/anchor distance in Mahalanobis units \\
\bottomrule
\end{tabular}
\end{center}

%======================================================================
\section{Common Experimental Setup}
%======================================================================

\subsection{Single-Class Protocol}

The scripts are designed to run one MVTec class at a time. This reduces
iteration time and makes failure analysis easier. The class is selected with
\code{CLASSNAME}; if unset, the default is \code{screw}.

\begin{lstlisting}[language=bash]
CLASSNAME=screw bash experiments/run_01_anchored_threshold.sh
\end{lstlisting}

\noindent
Common run-time overrides are:

\begin{lstlisting}[language=bash]
CLASSNAME=screw
DATAPATH=/path/to/mvtec
GPU=0
SEED=0
META_EPOCHS=40
GAN_EPOCHS=4
BATCH_SIZE=8
RESULTS_PATH=results
\end{lstlisting}

\subsection{Metrics}

The primary metrics are:

\begin{description}[leftmargin=1.2in, style=nextline]
  \item[Image AUROC]
  Measures whether the detector can rank defective images above normal images.

  \item[Pixel AUROC]
  Measures whether high anomaly scores land on anomalous pixels.

  \item[PRO]
  Region-level localization metric; useful when defects are spatial objects
  rather than isolated pixels.

  \item[Anomaly-pixel AUROC]
  Measures localization quality inside defective images. This is especially
  important because the current project state suggests image-level detection is
  ahead of pixel-level precision.
\end{description}

\subsection{Training Diagnostics}

The discriminator logs:

\begin{description}[leftmargin=1.2in, style=nextline]
  \item[\code{p\_true}]
  Fraction of real normal features classified past the positive margin.

  \item[\code{p\_fake}]
  Fraction of synthetic fake features classified past the negative margin.

  \item[\code{fake\_patch\_ratio}]
  Fraction of patches marked synthetic in sparse-mask experiments.

  \item[\code{loss}]
  Margin loss used to train the discriminator.
\end{description}

\noindent
The important warning is that a high \code{p\_fake} alone is not success. If the
fake samples are too easy or too far from realistic defects, the discriminator
can learn a boundary that improves training loss while hurting localization.

%======================================================================
\section{Experiment Ladder}
%======================================================================

\subsection{Experiment 0: Vanilla SimpleNet Noise}

\paragraph{Script.}
\begin{lstlisting}[language=bash]
bash experiments/run_00_simplenet_noise.sh
\end{lstlisting}

\paragraph{Configuration.}
\begin{lstlisting}[language=bash]
TSLRG_ANOMALY_MODE=simplenet_noise
TSLRG_PATCH_MASK_MODE=all
TSLRG_REFINE_STEPS=0
\end{lstlisting}

\paragraph{What it tests.}
This is the clean baseline. It restores SimpleNet-style feature noise inside
the same current training path~\cite{simplenet}:
\begin{equation}
  \mathbf{x}_{\mathrm{fake}}
    =
  \mathbf{x}_{\mathrm{real}} + \mathcal{N}(0, \sigma^2 I).
\end{equation}

\paragraph{Reasoning.}
Any proposed feature generator must beat this baseline under the same backbone,
projection, discriminator, loss, dataset split, and evaluation code. Without this
baseline, improvements may come from incidental training differences rather than
from the anomaly generator.

\paragraph{Expected behavior.}
This baseline should be stable and may localize surprisingly well because every
patch is fake and the perturbation is small. Its weakness is geometric: noise is
spent equally in every channel direction, including directions where normal
features have little or no meaningful variation.

\paragraph{Interpretation.}
If geometry-aware variants do not beat this baseline, then either the proposed
feature geometry is not aligned with real defects, or the synthetic perturbation
magnitude is miscalibrated. This experiment is the anchor for all claims.

%----------------------------------------------------------------------
\subsection{Experiment 1: Anchored PCA, Threshold Radius}

\paragraph{Script.}
\begin{lstlisting}[language=bash]
bash experiments/run_01_anchored_threshold.sh
\end{lstlisting}

\paragraph{Configuration.}
\begin{lstlisting}[language=bash]
TSLRG_ANOMALY_MODE=anchored
TSLRG_RADIUS_MODE=threshold
TSLRG_PATCH_MASK_MODE=all
TSLRG_REFINE_STEPS=0
\end{lstlisting}

\paragraph{What it tests.}
This is the previous best geometry-aware setting. Each fake feature is anchored
on a real normal feature and shifted inside the fitted patch PCA subspace:
\begin{equation}
  \mathbf{x}_{\mathrm{fake},p}
    =
  \mathbf{x}_{\mathrm{real},p}
    +
  r_p\,\Up\sqrt{\Lp}\,\mathbf{v}.
\end{equation}
The radius still follows the fitted Mahalanobis threshold:
\begin{equation}
  r_p = \sqrt{T_p} + \delta\,U(0,1).
\end{equation}

\paragraph{Reasoning.}
Anchoring removes one layer of modeling error: the base point is a real feature,
not a Gaussian draw. The PCA direction avoids wasting perturbation budget in
directions unsupported by normal data.

\paragraph{Expected behavior.}
This setting should preserve the earlier image-level gain. However, because
$\sqrt{T_p}$ can be around $40$ in a $1536$-dimensional feature space, the shift
may be much larger than subtle real defects. That can explain good image AUROC
but weak anomaly-pixel AUROC.

\paragraph{Interpretation.}
If image AUROC is strong but localization is weak, the issue is likely not the
anchor or subspace direction alone. The next suspect is radius magnitude.

%----------------------------------------------------------------------
\subsection{Experiment 2: Anchored PCA, Fixed Small Radius}

\paragraph{Script.}
\begin{lstlisting}[language=bash]
bash experiments/run_02_anchored_fixed_radius.sh
\end{lstlisting}

\paragraph{Configuration.}
\begin{lstlisting}[language=bash]
TSLRG_ANOMALY_MODE=anchored
TSLRG_RADIUS_MODE=threshold
TSLRG_RADIUS=1
TSLRG_PATCH_MASK_MODE=all
TSLRG_REFINE_STEPS=0
\end{lstlisting}

\paragraph{What it tests.}
This decouples anomaly magnitude from the fitted threshold:
\begin{equation}
  r_p = \rho\,U(0,1),
\end{equation}
where $\rho=\code{TSLRG\_RADIUS}$.

\paragraph{Reasoning.}
The statistical boundary $T_p$ answers a density-model question: how far is the
edge of the normal cloud in Mahalanobis units? But discriminator training may
need a different magnitude: small enough to mimic subtle defects, large enough
to be learnable.

\paragraph{Expected behavior.}
A fixed small radius may improve localization if the threshold-radius anomalies
were too extreme. It may reduce image AUROC if the anomalies become too subtle
for early discriminator training.

\paragraph{Interpretation.}
If this improves pixel metrics but hurts image AUROC, the generator needs a
curriculum or hard-negative schedule. If both improve, threshold radius was too
large for the task.

%----------------------------------------------------------------------
\subsection{Experiment 3: Patch-Calibrated Radius Sweep}

\paragraph{Script.}
\begin{lstlisting}[language=bash]
bash experiments/run_03_patch_radius_sweep.sh
\end{lstlisting}

\paragraph{Default sweep.}
\begin{lstlisting}[language=bash]
RADII="0.25 0.5 1 2 5"
\end{lstlisting}

\paragraph{Configuration.}
\begin{lstlisting}[language=bash]
TSLRG_ANOMALY_MODE=anchored
TSLRG_RADIUS_MODE=patch
TSLRG_PATCH_MASK_MODE=all
TSLRG_REFINE_STEPS=0
\end{lstlisting}

\paragraph{What it tests.}
The radius is normalized by the patch's fitted threshold and feature dimension.
This combines the patch-distribution view from PaDiM~\cite{padim} with the
distance-calibration motivation of CRAS~\cite{cras}:
\begin{equation}
  r_p = \rho \sqrt{\frac{T_p}{C}}\,U(0,1).
\end{equation}

\paragraph{Reasoning.}
Different patches can have different thresholds because their normal feature
variation differs. A global radius may under-perturb high-variance patches and
over-perturb low-variance patches. Patch normalization makes a radius sweep more
comparable across spatial locations.

\paragraph{Expected behavior.}
This should be more stable across classes and patches than a fixed radius. It is
especially useful if per-patch $T_p$ varies substantially.

\paragraph{Interpretation.}
If patch calibration beats fixed radius, then local scale matters. If it does
not, then the dominant issue may be direction distribution or sparse supervision
rather than per-patch magnitude.

%----------------------------------------------------------------------
\subsection{Experiment 4: Anchor-Gap Radius Sweep}

\paragraph{Script.}
\begin{lstlisting}[language=bash]
bash experiments/run_04_anchor_radius_sweep.sh
\end{lstlisting}

\paragraph{Default sweep.}
\begin{lstlisting}[language=bash]
RADII="0.25 0.5 1 2"
\end{lstlisting}

\paragraph{Configuration.}
\begin{lstlisting}[language=bash]
TSLRG_ANOMALY_MODE=anchored
TSLRG_RADIUS_MODE=anchor
TSLRG_PATCH_MASK_MODE=all
TSLRG_REFINE_STEPS=0
\end{lstlisting}

\paragraph{What it tests.}
The radius depends on how close the real anchor already is to the fitted
threshold, following the same high-level principle as distance-guided anomaly
synthesis in CRAS~\cite{cras}:
\begin{equation}
  r_{i,p}
  =
  \rho\,
  \max\left(
    \sqrt{T_p} - \sqrt{s(\mathbf{x}_{i,p})},
    0.05\sqrt{T_p}
  \right)
  U(0,1),
\end{equation}
where $s(\cdot)$ is the fitted Mahalanobis squared score.

\paragraph{Reasoning.}
Not all normal anchors are equally central. A feature already near the edge of
the normal distribution should require a smaller shift to become a useful
near-boundary negative. A very central feature can tolerate a larger shift.

\paragraph{Expected behavior.}
This should create more uniformly hard negatives: not trivially fake, but not
identical to normal features. It is the most direct test of the ``near-boundary
negative'' hypothesis.

\paragraph{Interpretation.}
If anchor-gap radius improves localization, the project has evidence that
negative sample difficulty matters. If it only improves image AUROC, sparse
patch supervision may still be needed.

%----------------------------------------------------------------------
\subsection{Experiment 5: Sparse Random Patch Sweep}

\paragraph{Script.}
\begin{lstlisting}[language=bash]
bash experiments/run_05_sparse_random_sweep.sh
\end{lstlisting}

\paragraph{Default sweep.}
\begin{lstlisting}[language=bash]
RATIOS="0.02 0.05 0.10 0.20"
\end{lstlisting}

\paragraph{Configuration.}
\begin{lstlisting}[language=bash]
TSLRG_ANOMALY_MODE=anchored
TSLRG_RADIUS_MODE=anchor
TSLRG_RADIUS=1
TSLRG_PATCH_MASK_MODE=random
\end{lstlisting}

\paragraph{What it tests.}
Only a random subset of patches is synthetic. Unselected patches are kept equal
to their real anchors, and the fake loss is applied only to selected patches.

\paragraph{Reasoning.}
Real industrial defects are usually sparse. Image-level synthetic anomaly methods
such as CutPaste~\cite{cutpaste}, NSA~\cite{nsa}, DRAEM~\cite{draem}, and
DeSTSeg~\cite{destseg} motivate using localized synthetic supervision. If every
patch in every fake image is
synthetic, the discriminator may learn image-level distribution shifts rather
than local defect evidence. Random sparse masks force the discriminator to
distinguish local deviations in the context of mostly normal features.

\paragraph{Expected behavior.}
Sparse random patches may improve pixel AUROC and PRO. Very small ratios may
under-train the fake class; very large ratios approach the all-fake setting.

\paragraph{Interpretation.}
If sparse random masks improve localization, the previous bottleneck was not
only feature magnitude; it was also the spatial training signal. If image AUROC
drops sharply, the fake signal may be too rare and needs a ratio curriculum.

%----------------------------------------------------------------------
\subsection{Experiment 6: Sparse Block Patch Sweep}

\paragraph{Script.}
\begin{lstlisting}[language=bash]
bash experiments/run_06_sparse_block_sweep.sh
\end{lstlisting}

\paragraph{Default sweep.}
\begin{lstlisting}[language=bash]
BLOCKS="3 5 7 9"
\end{lstlisting}

\paragraph{Configuration.}
\begin{lstlisting}[language=bash]
TSLRG_ANOMALY_MODE=anchored
TSLRG_RADIUS_MODE=anchor
TSLRG_RADIUS=1
TSLRG_PATCH_MASK_MODE=block
TSLRG_PATCH_MASK_RATIO=0.10
\end{lstlisting}

\paragraph{What it tests.}
Synthetic patches are selected in connected spatial blocks rather than
independent random positions.

\paragraph{Reasoning.}
Many defects have spatial continuity: scratches, stains, dents, missing material,
or texture corruption. Block masks are a crude feature-space analogue of
image-level synthetic anomaly masks. They may give a localization head a more
realistic spatial pattern than independent random patch corruption.

\paragraph{Expected behavior.}
Moderate block sizes should improve region-level metrics. Too-small blocks may
look like random isolated noise; too-large blocks may again become global
distribution shifts.

\paragraph{Interpretation.}
If block masks outperform random masks, spatial coherence is useful. If random
masks outperform blocks, real defects may be more local or the block sizes may be
too large for the chosen class.

%----------------------------------------------------------------------
\subsection{Experiment 7: Gradient-Guided PCA Refinement}

\paragraph{Script.}
\begin{lstlisting}[language=bash]
bash experiments/run_07_gradient_refinement_sweep.sh
\end{lstlisting}

\paragraph{Default sweep.}
\begin{lstlisting}[language=bash]
REFINE_STEPS_LIST="1 2 3"
REFINE_STEP_SIZES="0.02 0.05 0.1"
\end{lstlisting}

\paragraph{Configuration.}
\begin{lstlisting}[language=bash]
TSLRG_ANOMALY_MODE=anchored
TSLRG_RADIUS_MODE=anchor
TSLRG_RADIUS=1
TSLRG_PATCH_MASK_MODE=block
TSLRG_PATCH_MASK_RATIO=0.10
TSLRG_PATCH_MASK_BLOCK=5
\end{lstlisting}

\paragraph{What it tests.}
The generated fake features are refined using discriminator gradients before
the main discriminator update. The update is projected back into the fitted
patch PCA subspace and clamped by a Mahalanobis radius.

\paragraph{Reasoning.}
This experiment is motivated most directly by GLASS~\cite{glass}. Randomly
sampled negatives may be too easy or may miss the current decision
boundary. Gradient refinement converts the generator into a hard-negative
sampler: it searches for feature perturbations that the current discriminator
should learn to reject, while still respecting the local PCA geometry.

\paragraph{Expected behavior.}
One or two refinement steps may improve both image and pixel metrics by tightening
the decision boundary. Too many steps or too-large step sizes may create
adversarial artifacts that no longer resemble plausible defects.

\paragraph{Interpretation.}
If refinement helps, then the generator benefits from discriminator feedback.
If it hurts, the refinement may be producing adversarial directions that are
useful for fooling the discriminator but not aligned with real anomalies. The
next step would be to log before/after discriminator scores and feature norms.

%======================================================================
\section{Ablation Table Template}
%======================================================================

\begin{center}
\begin{tabular}{llllll}
\toprule
Experiment & Anchor & Radius & Mask & Refinement & Purpose \\
\midrule
0 & real & isotropic $\sigma$ & all & no & SimpleNet baseline \\
1 & real & threshold & all & no & previous anchored setup \\
2 & real & fixed & all & no & magnitude decoupling \\
3 & real & patch-calibrated & all & no & local scale calibration \\
4 & real & anchor-gap & all & no & near-boundary negatives \\
5 & real & anchor-gap & random & no & sparse local supervision \\
6 & real & anchor-gap & block & no & coherent local supervision \\
7 & real & anchor-gap & block & yes & hard negatives \\
\bottomrule
\end{tabular}
\end{center}

%======================================================================
\section{How to Interpret Outcomes}
%======================================================================

\subsection{Image AUROC Improves, Pixel Metrics Do Not}

The discriminator has learned image-level evidence, but the fake samples are not
teaching precise spatial localization. Prefer sparse masks, block masks, and
smaller or anchor-calibrated radii.

\subsection{Pixel Metrics Improve, Image AUROC Drops}

The synthetic signal may be too subtle or too sparse for robust image-level
ranking. Increase mask ratio, increase radius, or use a curriculum from stronger
to subtler anomalies.

\subsection{Training \code{p\_fake} Is High, Test Metrics Are Weak}

The fake anomalies are too easy or unrealistic. The discriminator is solving the
synthetic task rather than learning real-defect evidence. Reduce radius, use
anchor-gap calibration, or inspect PCA/residual energy.

\subsection{Gradient Refinement Hurts}

The refinement may be generating adversarial artifacts rather than plausible
defects. Reduce \code{TSLRG\_REFINE\_STEP\_SIZE}, reduce
\code{TSLRG\_REFINE\_STEPS}, or set \code{TSLRG\_REFINE\_MAX\_RADIUS}.

\subsection{Sparse Blocks Beat Sparse Random Masks}

Spatial coherence matters. This supports adding richer feature-space masks, such
as lines, ellipses, texture-like blobs, or masks copied from image-level anomaly
synthesis methods.

%======================================================================
\section{Recommended Execution Order}
%======================================================================

\begin{enumerate}[leftmargin=*]
  \item Run Experiment 0 and Experiment 1 on the same class. This establishes the
  vanilla-noise vs anchored-PCA comparison.

  \item Run Experiment 2, then the patch and anchor radius sweeps. Stop early if
  one radius family is clearly unstable.

  \item Take the best radius setting into sparse random and sparse block sweeps.

  \item Run gradient refinement only after a good sparse-mask configuration is
  found. Refinement is more expensive and easier to misread without a stable
  non-refined baseline.

  \item Repeat the winning configuration across additional classes.
\end{enumerate}

\noindent
For a quick smoke run:

\begin{lstlisting}[language=bash]
CLASSNAME=screw META_EPOCHS=1 GAN_EPOCHS=1 \
  bash experiments/run_01_anchored_threshold.sh
\end{lstlisting}

\noindent
For the full single-class suite:

\begin{lstlisting}[language=bash]
CLASSNAME=screw bash experiments/run_all_proposed_experiments.sh
\end{lstlisting}

%======================================================================
\section{Presentation Summary}
%======================================================================

\begin{itemize}[leftmargin=*]
  \item The baseline asks whether random feature noise is enough.
  \item Anchored PCA asks whether feature directions should follow the normal
  data manifold.
  \item Radius sweeps ask how close synthetic negatives should be to the normal
  boundary.
  \item Sparse masks ask whether localization requires local fake supervision.
  \item Gradient refinement asks whether the generator should adapt to the
  current discriminator.
\end{itemize}

\noindent
The strongest possible result would show a monotonic story: random noise
$\rightarrow$ anchored PCA improves image detection; calibrated radius improves
localization; sparse block masks improve PRO; gradient refinement improves the
hardest classes without destabilizing training.

\begin{thebibliography}{9}

\bibitem{simplenet}
Zhikang Liu, Yiming Zhou, Yuansheng Xu, and Zilei Wang.
\newblock SimpleNet: A Simple Network for Image Anomaly Detection and
Localization.
\newblock \emph{CVPR}, 2023.
\newblock \url{https://arxiv.org/abs/2303.15140}

\bibitem{padim}
Thomas Defard, Aleksandr Setkov, Angelique Loesch, and Romaric Audigier.
\newblock PaDiM: a Patch Distribution Modeling Framework for Anomaly Detection
and Localization.
\newblock \emph{ICPR}, 2021.
\newblock \url{https://arxiv.org/abs/2011.08785}

\bibitem{patchcore}
Karsten Roth, Latha Pemula, Joaquin Zepeda, Bernhard Schölkopf, Thomas Brox,
and Peter Gehler.
\newblock Towards Total Recall in Industrial Anomaly Detection.
\newblock \emph{CVPR}, 2022.
\newblock \url{https://arxiv.org/abs/2106.08265}

\bibitem{cutpaste}
Chun-Liang Li, Kihyuk Sohn, Jinsung Yoon, and Tomas Pfister.
\newblock CutPaste: Self-Supervised Learning for Anomaly Detection and
Localization.
\newblock \emph{CVPR}, 2021.
\newblock \url{https://arxiv.org/abs/2104.04015}

\bibitem{nsa}
Hannah M. Schl\"uter, Jeremy Tan, Benjamin Hou, and Bernhard Kainz.
\newblock Natural Synthetic Anomalies for Self-Supervised Anomaly Detection and
Localization.
\newblock \emph{ECCV}, 2022.
\newblock \url{https://arxiv.org/abs/2109.15222}

\bibitem{draem}
Vitjan Zavrtanik, Matej Kristan, and Danijel Sko\v{c}aj.
\newblock DRAEM: A Discriminatively Trained Reconstruction Embedding for Surface
Anomaly Detection.
\newblock \emph{ICCV}, 2021.
\newblock \url{https://arxiv.org/abs/2108.07610}

\bibitem{destseg}
Xuan Zhang, Shiyu Li, Xi Li, Ping Huang, Jiulong Shan, and Ting Chen.
\newblock DeSTSeg: Segmentation Guided Denoising Student-Teacher for Anomaly
Detection.
\newblock \emph{CVPR}, 2023.
\newblock \url{https://arxiv.org/abs/2211.11317}

\bibitem{glass}
Qiyu Chen, Huiyuan Luo, Chengkan Lv, and Zhengtao Zhang.
\newblock A Unified Anomaly Synthesis Strategy with Gradient Ascent for
Industrial Anomaly Detection and Localization.
\newblock arXiv, 2024.
\newblock \url{https://arxiv.org/abs/2407.09359}

\bibitem{cras}
Qiyu Chen, Huiyuan Luo, Haiming Yao, Wei Luo, Zhen Qu, Chengkan Lv, and
Zhengtao Zhang.
\newblock Center-aware Residual Anomaly Synthesis for Multi-class Industrial
Anomaly Detection.
\newblock arXiv, 2025.
\newblock \url{https://arxiv.org/abs/2505.17551}

\end{thebibliography}

\end{document}
